

To run the program using default settings:
\begin{lstlisting}[language=bash]
python main.py
\end{lstlisting}

\noindent To see all available arguments:
\begin{lstlisting}[language=bash]
python main.py -h
\end{lstlisting}

\noindent Descriptions on each argument are in the \hyperref[SEC:args]{following section}.

\noindent Following is the list of default arguments:\\
\begin{tabular}{r | l | l}
  \hline \hline
   & Value & Meaning \\
  \hline
  \texttt{-p} & ./data/VWT-data/circylinder & simulation directory \\
  \texttt{-m} &   & model name (\texttt{*.domain}) is automatically found \\
  \texttt{-A} & 0 180 181 & angular sweep settings: \texttt{np.linspace(0, 180, 181)} \\
  \texttt{-S} & 0 & phi-sweep \\
  \texttt{-F} & 100 300 & frequency band will be set to [100, 300] MHz \\
  \texttt{-v} &   & no validation \\
  \texttt{-a} & 0 & acquisition function: maximum variance \\
  \texttt{-x} & 1 & x-data normalization strategy: z-score (standard normalization) \\
  \texttt{-n} & 3 & number of initial frequency samples \\
  \texttt{--add} & 1 & added frequency samples per iteration \\
  \texttt{-t} & 3 & number of terms used for Gaussian process kernel \\
  \texttt{-s} & 1 & sampling strategy: Latin Hypercube Sampling \\
  \texttt{-i} & 20 & maximum iteration for Gaussian Process \\
  \texttt{-T} & 1e-5 & IVR tolerance; tolerance for Gaussian Process \\
  \hline \hline
\end{tabular}